% =============================================================================
% report.tex -- A soft FPGA on the PYNQ-Z2, reached over JTAG from a Pico
%
% Not compiled as part of any build. To produce a PDF:
%     pdflatex report.tex && pdflatex report.tex
% =============================================================================
\documentclass[11pt,a4paper]{article}

\usepackage[T1]{fontenc}
\usepackage[utf8]{inputenc}
\usepackage[margin=25mm]{geometry}
\usepackage{booktabs}
\usepackage{longtable}
\usepackage{listings}
\usepackage{xcolor}
\usepackage{hyperref}
\usepackage{graphicx}
\usepackage{amsmath}
\usepackage{enumitem}

\hypersetup{colorlinks=true, linkcolor=black, urlcolor=blue!60!black,
            citecolor=black, pdftitle={A soft FPGA on the PYNQ-Z2}}

\definecolor{codebg}{gray}{0.96}
\definecolor{codekw}{rgb}{0.13,0.29,0.53}
\definecolor{codecm}{rgb}{0.35,0.45,0.35}

\lstset{
  basicstyle=\ttfamily\scriptsize,
  backgroundcolor=\color{codebg},
  keywordstyle=\color{codekw}\bfseries,
  commentstyle=\color{codecm}\itshape,
  breaklines=true,
  frame=none,
  xleftmargin=8pt,
  framexleftmargin=8pt,
  showstringspaces=false,
  columns=fullflexible,
}
\lstdefinestyle{v}{language=Verilog}
\lstdefinestyle{py}{language=Python}
\lstdefinestyle{sh}{language=bash}

\newcommand{\file}[1]{\texttt{#1}}
\newcommand{\sig}[1]{\texttt{#1}}

\title{\textbf{A Soft FPGA on the PYNQ-Z2}\\[2mm]
       \large A 4$\times$4 CLB fabric, an IEEE 1149.1 TAP, and a bitstream
       compiler, driven from a Raspberry Pi Pico}
\author{Project report}
\date{}

\begin{document}
\maketitle

\begin{abstract}
\noindent
This project builds a miniature FPGA inside the programmable logic of a real
one. Sixteen configurable logic blocks sit in a routed mesh behind an
IEEE 1149.1 Test Access Port, configured by a 2896-bit scan chain and surrounded
by a boundary scan ring on the physical pins. A Python toolchain places, routes
and packs designs into that bitstream, and a Raspberry Pi Pico running DirtyJTAG
carries them to the board over four wires on a PMOD header.

Everything reported here has been verified twice: 174 checks in RTL simulation
and 163 input vectors replayed on the physical board, with expected results
produced by an independent software model of the fabric rather than by the
Verilog itself.
\end{abstract}

\tableofcontents
\newpage

% =============================================================================
\section{What was built, and in what order}

The work proceeded in three stages, each a complete and independently working
artefact. The identification code programmed into the TAP carries a version
nibble so the three can never be confused on the wire.

\begin{center}
\begin{tabular}{llll}
\toprule
\textbf{Stage} & \textbf{IDCODE} & \textbf{Top module} & \textbf{Config chain} \\
\midrule
1. Bare TAP           & \texttt{0x1BEEF093} & ---                    & none \\
2. Single CLB         & \texttt{0x2BEEF093} & \file{mini\_fpga\_top} & 71 bits \\
3. 4$\times$4 fabric  & \texttt{0x3BEEF093} & \file{fpga4x4\_top}    & 2896 bits \\
\bottomrule
\end{tabular}
\end{center}

Stage 1 established that a soft TAP written from scratch could be reached from
the Pico at all. Stage 2 added a logic block and a boundary scan ring and made
the board do something visible. Stage 3 scaled the block to a routed mesh and
added the toolchain that turns a design description into bits.

Stages 2 and 3 are both retained and both build from the same sources; the
Vivado script takes the desired top as an argument.

\subsection{The two JTAG chains}

A point of confusion worth settling immediately. Two entirely separate JTAG
chains are live at the same time and never interact.

\begin{center}
\begin{tabular}{llll}
\toprule
\textbf{Chain} & \textbf{Port} & \textbf{Contains} & \textbf{Used for} \\
\midrule
Built-in & micro-USB (PROG) & real \texttt{xc7z020} TAP + ARM DAP & loading the bitstream \\
Ours     & PMODA + Pico     & the soft TAP and the fabric          & everything afterwards \\
\bottomrule
\end{tabular}
\end{center}

The board's own JTAG is used exactly once per power cycle, to configure the
FPGA. From then on all interaction happens over the PMOD chain. The soft TAP is
not a pass-through or a tap onto the existing chain; it is a second, independent
IEEE 1149.1 device that happens to be implemented in the first one's fabric.

% =============================================================================
\section{Physical setup}

\subsection{Wiring}

Five wires between the Pico and PMODA. Both sides are 3.3\,V CMOS, so no level
shifting is required.

\begin{center}
\begin{tabular}{lllll}
\toprule
\textbf{Signal} & \textbf{Pico GPIO} & \textbf{Pico pin} & \textbf{PMODA pin} & \textbf{Zynq pin} \\
\midrule
TCK & GP18 & 24 & 7       & U18 \\
TMS & GP19 & 25 & 1       & Y18 \\
TDI & GP16 & 21 & 2       & Y19 \\
TDO & GP17 & 22 & 3       & Y16 \\
GND & GND  & 23 & 5 or 11 & --- \\
\bottomrule
\end{tabular}
\end{center}

Two details matter more than they appear to.

\textbf{TDI and TDO are not crossed.} JTAG signal names are chain-relative
rather than per-device: the probe's TDI pin is an output that feeds the target's
TDI input. Wiring them crossed like a UART produces a completely dead scan chain
with no error message.

\textbf{TCK is on PMODA pin 7, not pin 4.} TCK clocks every flip-flop in the
design, so it must reach a global clock buffer. \texttt{U18} is
\texttt{IO\_L12P\_T1\_MRCC\_34}, the only clock-capable pin on that header.
Placing TCK anywhere else raises a \texttt{CLOCK\_DEDICATED\_ROUTE} design rule
error. PMODB was rejected for exactly this reason despite being directly
connected rather than sitting behind PMODA's 200\,$\Omega$ series resistors,
which are electrically irrelevant below a few megahertz.

Both boards are powered from their own USB supplies with only ground connected
between them. Back-feeding one regulator from the other is a good way to damage
a board.

\subsection{The probe}

The Pico runs \emph{pico-dirtyJtag}, a port of Jean Thomas's DirtyJTAG to the
RP2040, presenting a vendor-specific USB interface at \texttt{1209:c0ca}. Both
\texttt{openFPGALoader} (cable \texttt{dirtyJtag}) and \texttt{urjtag} (cable
\texttt{DirtyJTAG}) speak it natively. Release V1.07 ships a prebuilt
\file{dirtyJtag.uf2}, so no SDK build is needed.

Firmware pin assignments, which the wiring table above matches:
TDI\,=\,GP16, TDO\,=\,GP17, TCK\,=\,GP18, TMS\,=\,GP19.

% =============================================================================
\section{The Test Access Port}

\file{rtl/jtag\_tap.v} is a complete IEEE 1149.1 TAP written from scratch. The
entire design lives in the TCK clock domain: there is no system clock, no Zynq
processing system, no AXI, and no reset pin. The probe clocks everything, which
is what a TAP is supposed to do and what removes every clock-domain-crossing
question from the design.

\subsection{Timing contract}

This is the part that is easy to get wrong by exactly one clock edge, and
getting it wrong produces a scan chain that looks almost right.

\begin{itemize}[nosep]
  \item TMS and TDI are sampled on the \textbf{rising} edge of TCK.
  \item The state machine advances on the \textbf{rising} edge.
  \item TDO is launched on the \textbf{falling} edge.
  \item IR and DR update latches fire on the \textbf{falling} edge.
  \item All shift registers are LSB-first.
\end{itemize}

\subsection{State encoding and a latent hazard}

The sixteen states use the conventional 4-bit JTAG numbering, chosen so that
\texttt{TEST\_LOGIC\_RESET} is \texttt{4'hF}. All-ones is also the natural
power-up value, which matters because there is no reset pin: FPGA flip-flops
come up from their bitstream INIT values, and the IEEE-guaranteed ``five TCKs
with TMS high'' sequence always returns to Test-Logic-Reset from anywhere.

The inherited \file{bsc\_cell.v} carries a header warning that Test-Logic-Reset
must be consumed \emph{synchronously}, because a combinational decode of the
state register can glitch. That warning was written against a different state
encoding, so it was checked against this one rather than assumed. A decode of
\texttt{4'hF} can glitch during any transition where source and destination
between them set every bit:

\begin{center}
\begin{tabular}{lll}
\toprule
\textbf{From} & \textbf{To} & \textbf{Bitwise OR} \\
\midrule
\texttt{RUN\_TEST\_IDLE 1100} & \texttt{SELECT\_DR 0111} & \texttt{1111} \\
\texttt{CAPTURE\_IR 1110}     & \texttt{EXIT1\_IR 1001}  & \texttt{1111} \\
\texttt{UPDATE\_IR 1101}      & \texttt{SELECT\_DR 0111} & \texttt{1111} \\
\bottomrule
\end{tabular}
\end{center}

Three such transitions exist. Nothing constrains four bits to change together,
so the warning applies here too and \sig{tlr} is only ever consumed on a clock
edge. The failure mode it prevents is unpleasant: the boundary silently commits
zeros at Update-DR, EXTEST and INTEST both fail, and capture and shift continue
to look perfect.

\subsection{Instruction set}

The instruction register is 4 bits. Capture-IR loads \texttt{4'b0001}, since
IEEE mandates the two least significant bits read back as \texttt{01} --- which
is exactly what lets a probe auto-discover the IR length with no database entry.

\begin{center}
\begin{tabular}{llll}
\toprule
\textbf{IR} & \textbf{Name} & \textbf{Register} & \textbf{Width} \\
\midrule
\texttt{0000} & EXTEST         & boundary & 9 \\
\texttt{0001} & SAMPLE/PRELOAD & boundary & 9 \\
\texttt{0010} & IDCODE         & idcode   & 32 \\
\texttt{0011} & USER           & user     & 32 \\
\texttt{0100} & INTEST         & boundary & 9 \\
\texttt{0101} & CONFIG         & config   & 71 or 2896 \\
\texttt{1111} & BYPASS         & bypass   & 1 \\
\bottomrule
\end{tabular}
\end{center}

All internal data registers shift in parallel and only the selected one reaches
TDO. This is cheaper than a shift-enable network and is safe because every
register re-captures at Capture-DR and only commits at an Update-DR qualified by
the instruction. The boundary strobes are qualified the same way, so scanning
IDCODE or CONFIG cannot disturb what the boundary is driving.

\subsection{The USER register}

Thirty-two bits. The low four are control, written at Update-DR; above them sit
\texttt{STATUS\_W} bits of live status, captured at Capture-DR.

\begin{center}
\begin{tabular}{lll}
\toprule
\textbf{Bit} & \textbf{Name} & \textbf{Effect} \\
\midrule
0 & \sig{ce}       & global clock enable, as a level \\
1 & \sig{sr}       & global synchronous set/reset, as a level \\
2 & \sig{cin}      & carry in to the bottom of each carry chain \\
3 & \sig{step}     & a 0$\rightarrow$1 transition gives exactly one clock enable pulse \\
4 & \sig{autostep} & each INTEST Update-DR gives one clock enable pulse \\
\bottomrule
\end{tabular}
\end{center}

\texttt{STATUS\_W} is a parameter: 3 for the single CLB, returning
\texttt{\{cout, o5, o\}}; 16 for the fabric, returning every tile's CLB output
in bits 19:4. That gives visibility inside the fabric without adding pins.

% =============================================================================
\section{The configurable logic block, from first principles}

The CLB, its LUT and the boundary cell were inherited unchanged from a prior
project and verified byte-identical with \texttt{cmp}. They are described here
in full because the bitstream format is defined by their internals.

\subsection{The fracturable LUT6}

\file{rtl/lut6.sv} is a 64-entry lookup table built as an explicit multiplexer
tree rather than an array index. The tree structure is the point: it makes the
\emph{fracture} tap fall out naturally.

\begin{lstlisting}[style=v]
wire [31:0] s0;  wire [15:0] s1;  wire [7:0] s2;  wire [3:0] s3;  wire [1:0] s4;

for (g = 0; g < 32; g++) assign s0[g] = i[0] ? init[2*g+1] : init[2*g];
for (g = 0; g < 16; g++) assign s1[g] = i[1] ? s0[2*g+1]   : s0[2*g];
for (g = 0; g <  8; g++) assign s2[g] = i[2] ? s1[2*g+1]   : s1[2*g];
for (g = 0; g <  4; g++) assign s3[g] = i[3] ? s2[2*g+1]   : s2[2*g];
for (g = 0; g <  2; g++) assign s4[g] = i[4] ? s3[2*g+1]   : s3[2*g];

assign o5 = s4[0];                  // LUT5 over INIT[31:0]
assign o6 = i[5] ? s4[1] : s4[0];   // the full 6-input function
\end{lstlisting}

Five levels of 2:1 multiplexers reduce 64 configuration bits to two values; the
sixth input selects between them. After five levels, \texttt{s4[0]} is already
the 5-input function of \texttt{i[4:0]} over \texttt{INIT[31:0]} --- a complete
LUT5, available for free. That is \sig{o5}, and it is what lets one physical
block implement two smaller independent functions.

An INIT word is simply the truth table: bit $n$ is the output for input vector
$n$. So \texttt{0x8888888888888888} is AND, because within each nibble only
index~3 ($i_1 i_0 = 11$) is set, and the nibble repeats so $i_{5..2}$ are
don't-care. \texttt{0x6666666666666666} is XOR by the same reasoning. The
toolchain builds these by evaluating a Python function over all 64 addresses
rather than by writing hexadecimal by hand.

\subsection{Carry chain and flip-flop}

\file{rtl/clb.sv} wraps the LUT with a carry chain and a flip-flop, mirroring
the Xilinx slice structure.

\begin{lstlisting}[style=v]
wire di      = cy_di_sel ? lut_o5 : i[0];
wire prop    = lut_o6;
wire cy_mux  = prop ? cin : di;     // MUXCY: propagate or generate
wire cy_sum  = prop ^ cin;          // XORCY: the sum bit
assign cout  = cy_en ? cy_mux : 1'b0;

wire comb    = cy_en ? cy_sum : (ff_d_sel ? lut_o5 : lut_o6);
wire ce_eff  = ff_ce_en ? ce : 1'b1;
wire sr_eff  = ff_sr_en ? sr : 1'b0;

always @(posedge clk)
  if      (sr_eff) q <= ff_rstval;  // synchronous, CE-independent
  else if (ce_eff) q <= comb;

assign o  = ff_en ? q : comb;
assign o5 = lut_o5;
\end{lstlisting}

In carry mode the LUT computes the \emph{propagate} term and the dedicated logic
does the rest: the sum is $\text{prop} \oplus \text{cin}$ and the carry out is
either \sig{cin} (propagate) or a generate term taken from \texttt{i[0]} or the
fractured LUT5. One block is therefore one bit of an adder, subtractor or
comparator.

The reset is synchronous and clock-enable-independent, matching FDRE/FDSE
behaviour: holding \sig{sr} keeps the flop reset, which is idempotent and
therefore safe to hold across a slow scan.

\subsection{The 71-bit CLB configuration word}

\begin{center}
\begin{tabular}{lll}
\toprule
\textbf{Bits} & \textbf{Field} & \textbf{Meaning} \\
\midrule
63:0 & \texttt{INIT}       & the LUT6 truth table \\
64   & \texttt{FF\_EN}     & \sig{o} is the flop's \sig{q} (1) or combinational (0) \\
65   & \texttt{FF\_RSTVAL} & value loaded on synchronous reset \\
66   & \texttt{FF\_CE\_EN} & the flop honours the global clock enable \\
67   & \texttt{FF\_SR\_EN} & the flop honours \sig{sr} \\
68   & \texttt{CY\_EN}     & carry-chain mode \\
69   & \texttt{CY\_DI\_SEL}& carry generate source: \texttt{i[0]} or \sig{o5} \\
70   & \texttt{FF\_D\_SEL} & datapath source: \sig{o6} or \sig{o5} \\
\bottomrule
\end{tabular}
\end{center}

% =============================================================================
\section{The tile, from first principles}

A CLB alone is not a fabric. What makes sixteen blocks into an FPGA is the
routing around them, and that is what a \emph{tile} adds. The architecture was
already specified in \file{rtl/clb\_pkg.sv}; the RTL implementing it is new.

\subsection{The shared source bus}

Every multiplexer in a tile --- and there are twenty-two of them --- chooses
from the same twenty-entry source bus. This uniformity is what keeps both the
hardware and the router simple.

\begin{center}
\begin{tabular}{lll}
\toprule
\textbf{Index} & \textbf{Source} & \textbf{Note} \\
\midrule
0     & constant 0            & \\
1     & constant 1            & \\
2     & this tile's CLB \sig{o}   & local feedback, no routing needed \\
3     & this tile's CLB \sig{o5}  & the fractured output \\
4--7  & tracks arriving from the North & \\
8--11 & tracks arriving from the East  & \\
12--15& tracks arriving from the South & \\
16--19& tracks arriving from the West  & \\
\bottomrule
\end{tabular}
\end{center}

Twenty sources need five select bits. Five bits address thirty-two, so twelve
encodings are unreachable; the hardware pads the bus to thirty-two entries with
hard zeros so a malformed bitstream reads a defined value rather than an
\texttt{X}.

That a tile can see its own outputs is deliberate and not merely a convenience:
it is what allows a single tile to hold a registered feedback loop --- a counter
bit, a toggle flip-flop --- with no routing resources at all.

\subsection{Connection box and switch box}

The two routing structures are the same circuit with a different multiplexer
count, so both are instances of \file{rtl/mux\_bank.v}.

\begin{center}
\begin{tabular}{llll}
\toprule
\textbf{Structure} & \textbf{Muxes} & \textbf{Bits} & \textbf{Drives} \\
\midrule
Connection box & 6  & 30 & the six LUT inputs \\
Switch box     & 16 & 80 & four outgoing tracks on each of four edges \\
\bottomrule
\end{tabular}
\end{center}

The connection box decides what the logic block \emph{sees}. The switch box
decides what the tile \emph{passes on}. Switch box multiplexer $j$ is indexed as
$j = \text{dir} \times \text{NTRACK} + \text{track}$ with
$\text{dir} \in \{N, E, S, W\} = \{0,1,2,3\}$.

\subsection{Tile configuration word}

\begin{center}
\begin{tabular}{lll}
\toprule
\textbf{Bits} & \textbf{Section} & \textbf{Width} \\
\midrule
70:0    & CLB            & 71 \\
100:71  & connection box & 30 \\
180:101 & switch box     & 80 \\
\midrule
        & \textbf{tile}  & \textbf{181} \\
\bottomrule
\end{tabular}
\end{center}

% =============================================================================
\section{The fabric}

\file{rtl/fabric.v} instantiates a $4 \times 4$ mesh. Coordinates increase North
and East; the tile index is $r \times 4 + c$, and that is also the order tiles
appear in the configuration word.

\subsection{Track adjacency}

Tracks are point to point between neighbours, and the track number is preserved
across the hop:

\begin{align*}
\text{tile}(r,c).\text{in\_n} &= \text{tile}(r{+}1,c).\text{out\_s} &
\text{tile}(r,c).\text{in\_s} &= \text{tile}(r{-}1,c).\text{out\_n} \\
\text{tile}(r,c).\text{in\_e} &= \text{tile}(r,c{+}1).\text{out\_w} &
\text{tile}(r,c).\text{in\_w} &= \text{tile}(r,c{-}1).\text{out\_e}
\end{align*}

Driving tile $T$'s switch box output in direction $D$ on track $t$ therefore
makes that signal arrive at $T$'s neighbour in direction $D$, on the opposite
side, on the same track. This single rule is the whole routing model, and the
router in the toolchain is a direct transcription of it.

The carry chain runs South to North up each column, with the bottom row's
\sig{cin} taken from the USER register, so a column of four tiles is a four-bit
adder.

\subsection{Where the pads meet the fabric}

Signals flow West to East, which makes a design read as a left-to-right
dataflow.

\begin{itemize}[nosep]
  \item \textbf{Inputs:} every row's West edge carries \texttt{pad\_i[3:0]} on
        tracks 0--3; every column's South edge carries \texttt{pad\_i[4]} on
        track 0 and \texttt{pad\_i[5]} on track 1.
  \item \textbf{Outputs:} \texttt{pad\_o[k]} is row $k$'s East edge, track 0,
        for $k = 0,1,2$.
\end{itemize}

Column 0 can therefore reach \texttt{pad\_i[3:0]} directly and row 0 can reach
all six; anything further in must be routed there through switch boxes. That
asymmetry is deliberate --- it is what makes this a fabric rather than a lookup
table with extra steps.

On the board, \texttt{pad\_i[5:0]} is \texttt{\{BTN3..BTN0, SW1, SW0\}} and
\texttt{pad\_o[2:0]} is \texttt{\{LD2, LD1, LD0\}}. LD3 indicates a non-zero
configuration word is loaded.

\subsection{Fabric configuration word}

\[
\text{FABRIC\_CFG\_W} = 181 \times 16 = 2896 \text{ bits}
\]

Tile 0 occupies bits 180:0, tile 1 bits 361:181, and so on. The chain shifts
LSB-first, so the first bit out of TDI is \texttt{INIT[0]} of tile $(0,0)$.

% =============================================================================
\section{The boundary scan ring}

Nine cells, all instances of the inherited \file{rtl/bsc\_cell.v}, a BC\_1
general-purpose cell with two flip-flops: a capture register on the scan path
and an update register that drives.

\begin{center}
\begin{tabular}{llll}
\toprule
\textbf{Bit} & \textbf{Cell} & \textbf{Captures} & \textbf{Drives} \\
\midrule
0   & \texttt{out\_o}    & CLB \sig{o}         & LD0 \\
1   & \texttt{out\_o5}   & CLB \sig{o5}        & LD1 \\
2   & \texttt{out\_cout} & CLB \sig{cout}      & LD2 \\
3--8& \texttt{in\_i0..5} & the six input pads  & the fabric's inputs \\
\bottomrule
\end{tabular}
\end{center}

Shifted LSB-first with the cell nearest TDO first, so the numbering is identical
whether reading or writing. Bits 8:3 are the vector applied; bits 2:0 are the
answers read.

The two-flop split is the point of the cell. The capture register samples the
system value at Capture-DR and then shifts, so moving data through the boundary
never disturbs what the cell is driving; the update register moves only at
Update-DR, so the boundary changes once, as a unit, at a moment the test
controls.

% =============================================================================
\section{The bitstream generation engine}

\file{host/bitstream.py} mirrors \file{rtl/clb\_pkg.sv} constant for constant
and adds placement, routing, packing and an independent software model.

\subsection{Interface}

\begin{lstlisting}[style=py]
from bitstream import Design, LUT

d = Design()
a, b   = d.input(0), d.input(1)
g_and  = d.lut(0, 0, LUT.and2(), [a, b])     # place an AND at tile (0,0)
g_nand = d.lut(2, 2, LUT.inv(0), [g_and])    # invert it three tiles away
d.output(0, g_and)
d.output(2, g_nand)
bs = d.build()                                # place, route, pack 2896 bits
\end{lstlisting}

\subsection{LUT INIT construction}

Truth tables are generated rather than written, by evaluating a Python predicate
over all 64 addresses:

\begin{lstlisting}[style=py]
def lut(fn, nin=6):
    init = 0
    for a in range(64):
        bits = [(a >> k) & 1 for k in range(6)]
        if fn(*bits[:nin]):
            init |= 1 << a
    return init
\end{lstlisting}

This produced \texttt{0x8888888888888888} for AND and
\texttt{0x6666666666666666} for XOR independently --- the same values that had
already been verified on hardware in stage 2, which was a useful early
confirmation that the engine and the silicon agreed about bit ordering.

\subsection{Routing}

The router is a breadth-first search over tiles. The start set is where the
signal is natively visible: for a cell, the tile producing it; for an input pad,
the edge tiles the pad arrives at. Each hop allocates one switch box
multiplexer.

\begin{lstlisting}[style=py]
for d in (DIR_E, DIR_N, DIR_S, DIR_W):
    nr, nc = r + DELTA[d][0], c + DELTA[d][1]
    for t in range(NTRACK):
        owner = self.sb_owner.get((r, c, d, t))
        if owner is not None and owner != key:
            continue                      # this track belongs to another net
        ncode = SRC_OF_DIR[OPPOSITE[d]](t)
        ...
\end{lstlisting}

A multiplexer already carrying \emph{the same} signal is free to traverse again.
That single condition is the whole of fanout support: a net reaching several
sinks reuses its existing path rather than allocating a second one, with no
special case anywhere.

Constants never route. A tile's own outputs never route. Both are already on
every tile's source bus.

\subsection{The software model}

\texttt{simulate(bs, pad\_i)} evaluates a packed bitstream in software and
returns \texttt{pad\_o}. It exists specifically so that expected results for the
testbench come from somewhere other than the Verilog.

It is an iterative fixed point over every CLB output and every track, not a
recursive walk. That is forced by the architecture: the mesh has cycles by
construction --- a tile can see its own output, and two tiles can point at each
other --- so there is no evaluation order to recurse in. The first version was
written recursively and immediately exhausted the stack, which is discussed in
Section~\ref{sec:defects}.

% =============================================================================
\section{Verification strategy}

The central idea is that the bitstream generator and the RTL must be made to
check each other, because a shared misunderstanding between them would otherwise
be invisible.

\begin{enumerate}[nosep]
  \item \file{sim/gen\_vectors.py} builds each design with the \emph{real}
        engine and computes expected outputs with the engine's \emph{software
        model}.
  \item It emits \file{sim/vectors.vh}: bitstream literals and expected
        input/output pairs.
  \item \file{sim/tb\_fpga4x4.v} loads the same bits into the RTL through the
        actual CONFIG scan chain and compares.
\end{enumerate}

If the generator and the fabric disagree about what a bitstream means, a named
vector fails. The two were written from the same specification but in different
languages, by different mechanisms, so agreement is meaningful.

The hardware self-test then replays the identical designs and vectors on the
board through INTEST, comparing against the same model.

\subsection{Results}

\begin{center}
\begin{tabular}{lll}
\toprule
\textbf{Stage} & \textbf{Coverage} & \textbf{Result} \\
\midrule
Fabric simulation      & 174 checks, 11 designs, 151 generated vectors & all pass \\
Single-CLB simulation  & 25 checks, 11 groups                          & all pass \\
Verilator lint         & both tops                                     & clean \\
Hardware self-test     & 11 designs, 163 vectors, on the board         & all pass \\
\bottomrule
\end{tabular}
\end{center}

Every configuration load on hardware completed in 0.04\,s with readback
verified.

% =============================================================================
\section{Defects found and fixed}
\label{sec:defects}

Recorded because most of them are instructive rather than clerical.

\subsection{Read-modify-write registers}

CONFIG, USER and the boundary are all capture/update registers: whatever is
shifted in is committed at Update-DR. A ``read'' that shifts zeros in returns
the correct answer and then leaves zeros behind it.

This was caught three separate times --- first in simulation when a USER
read-back silently cleared the register, then again when a CONFIG read-back
cleared the \sig{configured} flag. The fix throughout is to shift the same word
back in. \texttt{load\_bitstream()} does this and then verifies.

\subsection{The flip-flop could not be stepped from a USER scan}

An attempt to single-step the flop by writing the USER \sig{step} bit failed,
and the reason was structural rather than a coding error. \sig{bsr\_mode} is
decoded from the current instruction, as IEEE requires. The instant a USER scan
is loaded to issue the step, the boundary goes transparent and the CLB stops
seeing the vector INTEST applied --- it would clock in whatever the physical
switches happen to be.

The fix was to add \texttt{USER\_AUTOSTEP}. With it set, an INTEST Update-DR
raises the clock enable itself on the very next rising edge, while INTEST is
still selected and the boundary is still driving. One INTEST scan then means
``apply this vector and advance the flop one clock with it''.

\subsection{Recursive fabric evaluation}

The first software model evaluated a track by recursing into its source tile,
which recursed into its neighbours, and exhausted the stack immediately. Rewritten
as an iterative fixed point over all state, with an explicit oscillation bail-out.

\subsection{A stale TDO read}

\texttt{CMD\_GETSIG} returns the firmware's cached TDO, which is only refreshed
inside a PIO transfer. Reading it without first issuing a clock pulse returns a
stale bit. All host tools now pulse TCK and then read.

\subsection{Vivado combinational loop DRC}

Implementation failed with:

\begin{lstlisting}
[DRC LUTLP-1] Combinatorial Loop Alert: 1569 LUT cells form a combinatorial loop.
\end{lstlisting}

This is the architecture, not a defect. A routing mesh in which any switch box
multiplexer can select any neighbouring track has cycles in the netlist graph by
construction; whether a loop actually \emph{exists} depends on the bitstream,
and none the router emits contains one.

Of Vivado's two documented resolutions, naming a net per loop with
\texttt{ALLOW\_COMBINATORIAL\_LOOPS} is unworkable here: there are many loops and
the net names are synthesis output that changes every build. The check is
therefore downgraded, in \file{constr/pynq\_z2\_jtag.xdc} and re-applied by
\file{vivado/drc\_waiver.tcl} as a hook on \texttt{opt\_design} and
\texttt{write\_bitstream}, because a DRC severity is a tool setting rather than a
design property and does not reliably survive every step.

The downgrade is scoped to \file{fpga4x4\_top}. The single-CLB build has no mesh
and must never trip LUTLP-1; if it does, that is a genuine defect and masking it
would be wrong.

\subsection{An empty worst-slack value read as a failure}

The build reported \texttt{timing not met (WNS  ns)} --- with no number. Vivado
had returned an empty slack value, and the Tcl guard \texttt{if \{\$wns < 0\}}
falls back to \emph{string} comparison when the operand is not numeric, where
\texttt{""} sorts before \texttt{"0"}. An absent measurement was therefore
reported as a failure. The check now confirms \texttt{string is double -strict}
before comparing.

An empty value is expected on this design rather than alarming: with loops
broken arbitrarily for analysis there may be no complete path for the tool to
call the worst one.

\subsection{Terminal echo in the interactive console}

Only the first keystroke registered. Two causes. \texttt{tty.setcbreak} disables
line buffering but \emph{leaves ECHO on}, so every key was echoed into the middle
of a display that redraws with cursor-up. And \texttt{select} watched the raw
file descriptor while \texttt{sys.stdin.read(1)} read through Python's buffered
text wrapper, stranding keys in the buffer. Fixed by clearing \texttt{ECHO}
explicitly while keeping \texttt{ISIG}, and by reading with \texttt{os.read} on
the descriptor.

\subsection{LEDs dark during INTEST}

The scan reported \texttt{LEDs = 001} while the board sat unlit. Both were
correct. In INTEST \emph{every} cell drives, including the three output cells, so
the pads are fed from those cells' update registers rather than from the fabric.
Shifting zeros into bits 2:0 actively held the LEDs off; the fabric's answer
arrives only by capture and never reaches a pad on its own.

The console now uses three scans per step: apply the vector, capture the answer,
then write that answer back into the output cells. The LEDs mirror the fabric one
scan behind. A negative control --- running the previous two-scan code against a
faithful model of \file{bsc\_cell} --- reproduced the original symptom exactly,
confirming the test detects the defect rather than merely agreeing with the fix.

\subsection{Loading 2896 bits}

One USB round trip per bit is acceptable for a 32-bit IDCODE and hopeless for the
configuration chain. \texttt{shift\_dr\_fast()} uses DirtyJTAG's
\texttt{CMD\_XFER}, which shifts up to 496 bits per packet with TMS held low.

Two details had to be right. Data is MSB-first \emph{within each byte}, pinned
down by the firmware's \texttt{jtag\_set\_tdi} writing \texttt{1u << 7} for a
single bit. And only whole bytes go that way; the ragged tail and the final
TMS-high bit are ordinary pulses. Every load is verified by readback and falls
back to the per-pulse path on disagreement. Measured: 0.04\,s per load.

% =============================================================================
\section{File manifest}

\subsection{RTL}

\begin{longtable}{llp{72mm}}
\toprule
\textbf{File} & \textbf{Lines} & \textbf{Purpose} \\
\midrule
\endhead
\file{clb\_pkg.sv}     & 237 & Architecture definition: field offsets, widths, source encodings, LUT builders. \emph{Inherited, byte-identical.} The authority for the bitstream format. \\
\file{lut6.sv}         & 35  & Fracturable 6-input LUT as a mux tree. \emph{Inherited, byte-identical.} \\
\file{clb.sv}          & 63  & LUT + carry chain + flip-flop. \emph{Inherited, byte-identical.} \\
\file{bsc\_cell.v}     & 122 & BC\_1 boundary scan cell. \emph{Inherited, byte-identical.} \\
\file{jtag\_tap.v}     & 283 & IEEE 1149.1 TAP: state machine, IR, IDCODE/USER/CONFIG/BYPASS, boundary strobes, step logic. \\
\file{mux\_bank.v}     & 41  & $N$ configurable multiplexers on a shared source bus; used as both connection box and switch box. \\
\file{tile.v}          & 102 & One tile: CLB + connection box + switch box, and the 20-entry source bus. \\
\file{fabric.v}        & 144 & The $4\times4$ mesh: neighbour wiring, edge ports, carry chains. \\
\file{fpga4x4.v}       & 187 & Fabric core: TAP + boundary ring + fabric, and the pad-to-edge mapping. \\
\file{fpga4x4\_top.v}  & 56  & Board wrapper: BUFG on TCK, switches, buttons, LEDs. \\
\file{mini\_fpga.v}    & 153 & Single-CLB core (stage 2), retained as a fallback. \\
\file{mini\_fpga\_top.v}& 69 & Board wrapper for the single-CLB build. \\
\bottomrule
\end{longtable}

\subsection{Host toolchain}

\begin{longtable}{llp{72mm}}
\toprule
\textbf{File} & \textbf{Lines} & \textbf{Purpose} \\
\midrule
\endhead
\file{bitstream.py}     & 489 & Architecture model, LUT builders, placement, router, packer, software simulator. \\
\file{designs.py}       & 118 & The eleven example designs, shared by simulation and hardware so both run the same thing. \\
\file{dirtyjtag.py}     & 342 & Shared transport: USB, TAP navigation, per-pulse and bulk shifting, boundary helpers. \\
\file{fpga.py}          & 242 & Fabric tool: report, load, self-test, watch, boundary-scan console. \\
\file{bscan.py}         & 202 & Interactive boundary-scan control: INTEST/EXTEST/SAMPLE. \\
\file{minifpga.py}      & 345 & The equivalent tool for the single-CLB build. \\
\file{tap\_probe.py}    & 144 & Low-level diagnostic: raw TDO stream, loopback, LED walk. \\
\file{wire\_check.py}   & 179 & Per-jumper continuity and PMODA pin-numbering check. \\
\file{detect.sh}        & 43  & IDCODE read via openFPGALoader and urjtag; names the build found. \\
\file{install\_urjtag\_part.sh} & 105 & Adds both builds to urjtag's part database. \\
\bottomrule
\end{longtable}

\subsection{Simulation, constraints and build}

\begin{longtable}{llp{72mm}}
\toprule
\textbf{File} & \textbf{Lines} & \textbf{Purpose} \\
\midrule
\endhead
\file{tb\_fpga4x4.v}      & 261 & Fabric testbench; includes the generated vectors. \\
\file{tb\_mini\_fpga.v}   & 338 & Single-CLB testbench, eleven groups. \\
\file{gen\_vectors.py}    & 52  & Turns the designs into \file{vectors.vh}. \\
\file{vectors.vh}         & 235 & \emph{Generated.} Bitstream literals and expected outputs. \\
\file{run\_fabric\_sim.sh}& 19  & Regenerate vectors, compile, run. \\
\file{run\_sim.sh}        & 18  & Single-CLB simulation. \\
\file{lint.sh}            & 33  & Verilator lint of both tops, with documented waivers. \\
\file{pynq\_z2\_jtag.xdc} & 117 & Pins, pulls, clock, timing, DRC handling; serves both tops. \\
\file{create\_project.tcl}& 229 & The entire Vivado flow in one session: create, build, optionally program. \\
\file{drc\_waiver.tcl}    & 23  & Re-applies the LUTLP-1 downgrade at the steps where it runs. \\
\bottomrule
\end{longtable}

% =============================================================================
\section{Operation}

\subsection{Build}

\begin{lstlisting}[style=sh]
vivado -mode batch -source vivado/create_project.tcl                  # build
vivado -mode batch -source vivado/create_project.tcl -tclargs all     # build + program
vivado -mode batch -source vivado/create_project.tcl -tclargs program # program only
\end{lstlisting}

Adding a top module name builds the single-CLB design instead. The whole flow
runs in one Vivado session; splitting create and build across two scripts meant
the second had to re-open a project the first had already opened.

\subsection{Use}

\begin{lstlisting}[style=sh]
./sim/run_fabric_sim.sh              # 174 checks
./host/fpga.py --list                # the built-in designs
./host/fpga.py --report three        # place and route, no hardware needed
./host/fpga.py --load showcase       # compile, load, run, watch the switches
./host/fpga.py --selftest            # every design, verified on the board
./host/fpga.py --load xor --bscan    # take the switches over from JTAG
\end{lstlisting}

% =============================================================================
\section{Known limitations}

\begin{itemize}
  \item \textbf{Static timing analysis is inconclusive.} The fabric's
        combinational loops are broken arbitrarily by the tool, so some paths are
        not analysed and any reported worst slack is indicative rather than
        exhaustive. TCK is constrained at 1\,MHz for the fabric against a link
        driven at 100\,kHz--1\,MHz; the margin is deliberate compensation.
  \item \textbf{The router is shortest-path, not congestion-aware.} It allocates
        the first free track on a breadth-first path. For designs of the size the
        fabric holds this has been sufficient, but it will fail on a congested
        design where a longer route would have succeeded.
  \item \textbf{The software model is combinational only.} It rejects designs
        with \texttt{FF\_EN} set. Sequential designs can be built and run on
        hardware but their expected results are not independently generated.
  \item \textbf{No timing model.} The toolchain has no notion of delay, so it
        cannot report a maximum frequency for a placed design.
  \item \textbf{Three output pads.} The boundary ring exposes three outputs and
        six inputs, which bounds what a design can show without additional pins.
\end{itemize}

\end{document}
